\chapter{客观精神与实践哲学}

客观精神学说为黑格尔解读提出了一系列难题，因此对属于客观精神的法哲学做出更进一步的阐释之前，任何解读都必须首先应对这些难题。我们可以将其归为三个问题域：(1) 这一概念的形成问题；(2) 哲学体系中的安排问题；(3) 特别难以处理的问题，即黑格尔思想发展的不同时期对这一概念的处理问题。第一个问题相应于这一概念自身所包含的多义性问题。我在这里只强调在黑格尔的语言使用中交替出现的几个最重要的意义成分。客观精神是超个体、跨主体的精神，即两个或更多的主体或者说主体统一体（家庭、等级、团体、民族等）之共同物。在这种意义上，客观精神就是普遍精神（《哲学全书》第1版，第399节）。依照普遍意志的模式——这种普遍意志并不是从个体意志中归纳出来的，而是贯穿于这种个体意志之中，并且统摄这些个体意志的——客观精神将个体的、“主观的”精神包含于自身之中。不同于普遍意志的公式——这种公式表达了一种抽离了时间和生成的“内在”意志关系——，普遍精神作为客观精神处于外部现象之中。就此而言，对于黑格尔来说，客观精神本质上就是历史精神。而外在显现，即客观精神的对象性或者现实性，又包含了许多方面，其中概念的意义繁多，模棱两可，时而指对象性的或者现实的精神，时而指对象化的或者现实化的精神，最后还指对象化了的或者现实化了的精神（或者用对这种模棱两可不满的尼古拉·哈特曼的话来说，就是“客观”精神和“客观化了的”精神）\footnote{《精神存在的问题》，第2版，柏林，1949年，第196页以下。}。

在涉及黑格尔思想中的体系安排及其发展演变的问题时，我希望自己限于如下主张。客观精神概念只是在1817年《哲学全书》中才与众所周知的主观精神、客观精神以及绝对精神的划分一道引入。早期作品中没有这种划分。尽管青年黑格尔的种种体系安排遵循了由笛卡尔和康德所规定的现代哲学的理性、自然与精神的三一体模式，但是精神哲学还没有在自身内部进行这样的区分，而是最初冠以“伦理体系”和“实践哲学”这样的名称。此时“精神”还完全是“绝对的”，以至于黑格尔所知的唯一区分涉及的是“绝对精神”自身：在一族人民中，精神是“存在着的绝对精神”\footnote{参见《耶拿实在哲学》第2卷（Jenenser Realphilosophie II），1805—1806年，J. 荷夫迈斯特出版，莱比锡，1931年，第272页。}。直到《纽伦堡哲学入门》之《哲学全书》（1811—1812年）中才出现了区分，尽管此时黑格尔似乎尚未拥有后来那样的成熟概念。这里，黑格尔区分了精神的实在化与其完成两个领域（第128节），但是前者并不是在“客观精神”而是在“实践精神”的标题下（第173节）加以探讨的\footnote{参见黑格尔：《全集》（Werke）第3卷，H. 格洛克纳出版，斯图加特，1927年，第200、215页。引文来自格洛克纳出版的《黑格尔全集》（百年纪念版），斯图加特，1927年之后，以下简称《全集》第1卷，等等。}。对于客观精神学说的概念形成与体系安排来说，这具有决定性的意义。《纽伦堡哲学入门》草稿仍然在概念上坚持了一种为黑格尔以及受其支配的19世纪哲学史研究所逐渐放弃的见解：“客观精神”的起源及其与传统称之为“实践哲学”的那个哲学部门之间的联系。要解释这种联系，黑格尔精神哲学中用来对概念进行处理的那些观点是不够的。当然我们可以说，这些观点对于评价实践哲学在近代思想中所起的作用来说，完全是致命的。[其后，]一直到《纽伦堡哲学入门》都没有分开的客观精神与实践精神开始分道扬镳。在1817年百科全书体系中，实践精神成为了主观精神的一个组成部分，即心理学的意志理论（第368-399节）。毫无疑问，实践哲学陷于遗忘根本上要归因于精神科学中的心理学理论，这与黑格尔体系中实践精神消融于主观精神以及自然哲学与精神哲学的对立有关。受这种对立的影响，狄尔泰创立了他的学派。他试图证明，早在17世纪，随着意识分析以及人类学情感理论的发现，这样一种对立就已经开始取代理论哲学与实践哲学的“错误划分”\footnote{参见《精神科学引论》（Einleitung in die Geisteswissenschaften）第1卷（1883年），收入《全集》第1卷，第4版，斯图加特/哥廷根，1959年，第225、378页。}。由此，狄尔泰基于精神哲学的体系构想对青年黑格尔的作品所做的解释也就时常误入歧途。在这种情况下，就完全不会提出“黑格尔用客观精神学说取代实践哲学到底意味着什么”这一问题。而且同样也不会问，这种概念的转变在历史上究竟意味着什么，以及为什么黑格尔只有在其哲学发展相对较晚的时候才会形成这个概念。为了对这些问题进行回答，我们就必须反其道而行之：从重构青年黑格尔的实践哲学观念出发，到其消融于“客观精神”学说以及作为其核心部分的法哲学之中。我首先尝试把黑格尔置入现代思想与传统之争的思想脉络之中，以便确定黑格尔得以介入这一争论过程的确切位置。

\section*{1}

由亚里士多德所创立的实践哲学学科及其中世纪和近代早期传统的突出之处，恐怕不是其理论内容，而是其在这些世纪中所拥有的独特形式和历史效力。黑格尔在《哲学史讲演录》中已经指出了这一点，并且指出了产生这一效力的原因。他说：“一般说来，一直到最近，我们所见到的还不是真正的思辨性的实践哲学。”\footnote{《哲学史讲演录》第1卷，C. L. 米希勒出版，《黑格尔全集》第17卷，第291页。（中译文引自〔德〕黑格尔：《哲学史讲演录》第一卷，贺麟、王太庆译，商务印书馆1960年版，第248页。——译者）}实践哲学之所以恒定不变，是由于其非思辨的特征。即是说，实践哲学的诸多公理独立于“第一哲学”的前提。对亚里士多德及其支配下的经院哲学来说，尽管实践哲学运用了形而上学的公理，但是两者之间并不存在直接的奠基关系。随着黑格尔所言的那种现代的到来，也就是17和18世纪自然法的发展，这一情况发生了改变。此时理论延伸到实践哲学的结构之中，实践哲学也提出了那种为现代科学方法论理想所推崇的普遍论证的要求。对于霍布斯而言，实践哲学的任务在于，如同认识图形的大小比例般，以同样的精确性认识人类行为的诸种关系。如同几何学一样，实践哲学现在也应当成为一门先天的、可以演证的科学。在霍布斯看来，实践哲学的对象非常适合于这一目的，因为对人来说，只有人自己产生的那些事物，才能有一种先天的、出于“第一根据”（ex causis）的科学\footnote{霍布斯：《论人》（De homine），第10章，第5节。}。

现代实践哲学的“思辨”特征在于这种对几何学认识理想的认同，而这同时也是批判实践哲学传统的第一步。第二步则是从实践哲学自身的经院形态出发，这种经院形态在18世纪初将它从外部通过先天演绎的方法论要求施加给它的批判原则纳入自身之中。沃尔夫对这一方法的运用导致实践哲学的体系分裂为一个“理性的”部分和一种“经验的”部分。理性部分对一般人类行动的区分以及权利和义务原则进行先天的证明，并且模仿笛卡尔的普遍数学（mathesis universalis）的观念，冠以普遍实践哲学（philosophia practica universalis）\footnote{《按照数学方法撰写的普遍实践哲学》（Philosophia practica universalis mathematico methodo conscripta），莱比锡，1703年。那篇莱比锡辩论论文——沃尔夫说是他写的，因为笛卡尔未曾研究过这一领域[《讲授纲要》（Ratio praelectionum），第6章，第3节]——构成了《普遍实践哲学》（第1-2卷，法兰克福/莱比锡，1738-1739年）的体系基础。}之名。这一新科学的引入为近代思想与实践哲学传统之争的第三步做好了准备。这一步由康德迈出。康德在“著名的沃尔夫之道德哲学导论”中清楚地意识到，这门新科学距离实践理性的批判有多远。这门新科学的引入源于对来自传统的概念和对象进行论证的需要，而普遍实践哲学则在方法和内容上仍然保持了与这些概念和对象的联系。依照康德的观点，这也是它没有进入一片“新领域”的原因之所在，因为它作为普遍的实践哲学仅仅把意志的普遍概念、“一般意愿”作为线索，因而并不考察一种“特殊”意志的可能，“比如一种无须任何经验性的动因完全为先天原则所规定，人们可以称之为纯粹意志的意志”\footnote{《道德形而上学奠基》（Grundlegung zur Metaphysik der Sitten，1785年，科学院版），第4卷，第390页。（中译文引自〔德〕康德：《道德形而上学奠基》，杨云飞译，邓晓芒校，人民出版社2013年版，第6页。——译者）}。

这种“特殊的”、从所有自然联系中摆脱出来的意志的发现，使得康德用实践哲学自身论证的观念取代了对实践哲学“特殊部分”（即传统的经院形态）的论证。同时，这种自身论证的观念也使得对于沃尔夫的普遍实践哲学来说具有构建意义的向传统流传下来的概念和对象的回溯成为了不可能的事情。康德《实践理性批判》（1788年）的主题在于，实践知识的最高条件不再像沃尔夫的“导论”所言的那样从自然规律中，而是从那个纯粹意志的理念中演绎出来，此纯粹意志自身在其自由中即是法则。自由的法则摧毁了实践哲学传统结构的一贯性，以及由沃尔夫再次建立起来的传统实践哲学的“普遍”部分与“特殊”部分之间的联系。政治学、家政学和伦理学作为明智和幸福的理论，完全被排除在实践哲学之外；对康德来说，它们都只是实用科学和技术性科学，属于经验性的经验世界、明智和技巧的活动领域。它们可能会从自然联系中借来规则，以便每次都能产生出按照因果律来说是可能的结果；然而它们并不能从自然联系中借来那些——如在旧的明智理论中——赋予伦理和政治行动以伦理义务的力量的目的。康德的自由法则终结了这种借用。依照康德的观点，只要它包含了所有行动的先天规定根据，那它就是一种“被称为实践哲学的特殊哲学”的对象，从而与那些违法地占据了“实践”之名义的（康德意义上的）实用学科和技术性学科相对立\footnote{《判断力批判》，导言。《全集》（科学院版）第5卷，第171页以下。}。

黑格尔在康德自由概念及其理论前提（即《纯粹理性批判》中发现的自我意识的自发性）的基础上，对实践哲学的经院传统进行了批判。这一批判在最大程度上超越了这个传统。辩证法完成了几何学方法、理性方法和先验方法未能完成的事情：传统的瓦解。这并不是因为黑格尔思想的尖锐批判和激进（在这个方面，霍布斯和康德都要明显胜过他），而是因为它们得以产生的条件。我们已经指出第一个条件：康德的意志自律的思想。黑格尔赞同这一思想，这使得他对传统的批判从一开始就超越了霍布斯和沃尔夫的批判。因为在霍布斯和沃尔夫那里，自由问题以及亚里士多德传统中与自由问题相关的主奴关系并没有得到解决。第二个条件是把现代国民经济学纳入实践哲学的构建中，因此也就是尝试将实践哲学的传统原则嵌入社会和工业世界的历史当下经验之中。由此产生的最重要后果就是黑格尔对制造（$\pi o\acute{\iota}\eta\sigma\iota\varsigma$）与行动（$\pi\rho\tilde{\alpha}\xi\iota\varsigma$）之间关系的思辨瓦解。尽管霍布斯已经将两者间的传统界线相对化了，但还没有真正克服，因为这种相对化还局限于政治领域，也就是人造的作品“利维坦”之内，还没有进入经济领域。在康德那里，这种情况再次出现。康德试图把当时刚刚确立的“国家经济学”从“特殊的”实践哲学中完全排除掉。最后，第三个条件，尽管它已经是黑格尔思想的成就，而不只是一个前提，但它涉及与历史维度的联系。这种联系使得实践哲学的范畴从完全无历史的经院传统中摆脱出来，使它们进入——将自己带入一种与这些范畴的内容完全改变了的关系之中的——“概念”的那种辩证运动之中。

同时，由黑格尔所引起的实践哲学史上的转折是从一个悖论开始的。如果我们想要理解这种转折的影响，那就必须认识并且说出这一悖论。这个悖论就是，一方面，黑格尔与传统保持了最遥远的距离，同时却又最先重新发现了一条直接通往这一传统的通道。这一点将他与传统的争论同霍布斯、沃尔夫和康德等人对传统的不同批判区别开来。就他们——如沃尔夫和康德——与亚里士多德主义经院哲学传统的一种特定的、相对较晚的阶段发生联系而言，他们与亚里士多德要么处于一种历史上相距遥远的关系中（霍布斯），要么处于一种仅仅是间接的关系中，青年黑格尔则研究了亚里士多德伦理学和政治学\footnote{黑格尔在斯图加特文理中学时期就开始研读《尼各马可伦理学》。参见罗森克兰茨（K. Rosenkranz）：《黑格尔传》（Hegels Leben），柏林，1844年，第11页。黑格尔对亚里士多德的真正研究（借助于16世纪的巴塞尔老版）则要追溯到在图宾根神学院的那些岁月。参见D. 亨利希：《洛伊特魏因论黑格尔：一份黑格尔传记文献》（D. Henrich, Leutwein über Hegel. Ein Dokument zu Hegels Biographie）一文中洛伊特魏因的证明，见《黑格尔研究》（Hegel-Studien）第3卷，1965年，第58页。《政治学》阅读的影响反复出现在耶拿时期的早期作品中。参见《全集》第1卷，第111、114、485页以下、494页以下、511页。《政治学与法哲学文集》（Schriften zur Politik und Rechtsphilosophie），G. 拉松出版，莱比锡，1913年，第417页以下、442页以下、460页以下、478页以下。当指亚里士多德的政治学和伦理学。——译者}，由此作为德国最早的人士之一，对传统获得了一副比那些批判者要恰当得多的形象，[反之，]对[霍布斯、沃尔夫和康德]这些批判者来说，这些渊源要么完全不存在，要么极端模糊不清。但是，这一过程中至关重要的倒不是黑格尔大体上重构了古典理论，而是他试图把古典理论同时运用到自己的思想以及当下历史经验向实践哲学所提出的任务上。黑格尔耶拿时期的著作即基于这种立场。他要求应用古典传统，从而不得不将一大堆材料放到一些概念之下，这些概念通常并不是来自这些材料，而是有着相反的来源。因而在黑格尔处理实践哲学的第一阶段（1801—1804年），概念与实在、形式与内容仍然是完全分裂的。尽管在第二阶段（1805—1807年）他逐渐对自己的哲学基础及其与实践领域的关系有了清晰的理解，但是只有在第三阶段——我们可以把两卷本的《逻辑学》以及《哲学全书》的出版（1816—1817年）作为这一阶段的开端——他才真正形成了“客观精神”这一概念。随着这一概念的形成，他克服了实践哲学传统的原则和体系形式。

\section*{2}

黑格尔在耶拿时期之初回归古典政治哲学并非什么意外之事。在18世纪90年代，黑格尔主要集中于神学和历史研究时就已经隐隐约约展现了这种倾向，但其意义有所不同。黑格尔早期作品的吊诡并不在于一般而言他又重新找到了回归古典城邦伦理的道路，而在于他同时按照康德《实践理性批判》的概念（自由、自律、自发性等）来对它进行衡量。在他看来，依据实践理性而行动和信仰的个体与古代的共和主义者都是自由的；前者如他的“精神性的、不朽的”理性本质所要求的那样来规定自己，后者则“在其人民的精神中”履行他的义务\footnote{《早期神学著作集》（Theologische Jugendschriften），H. 诺尔出版，图宾根，1907年，第70页。参见埃弗拉伊姆（F. Ephraim）：《黑格尔早期作品中的自由概念研究》（Untersuchungen über den Freiheitsbegriff Hegels in seinen Jugendarbeiten），见《哲学研究》（Philosophische Forschungen）第7卷，柏林，1928年，第59、63页。近著参见歌兰德（I. Görland）：《青年黑格尔的康德批判》（Die Kant-Kritik des jungen Hegel），美因河畔法兰克福，1966年，第4页以下。}。康德《道德形而上学》（1797年）的出版结束了这种混杂，因为康德的这部著作并没有为自律思想向政治世界这样一种拓展提供什么凭靠，并且遭到了法兰克福时期黑格尔（1798—1799年）的批评。黑格尔从古代实体性伦理、从一种人民宗教的完美自由中理性和自然的同一（Einssein）出发；但是古代的自由既没有构成他自己政治意愿的标准，他也没有用来解释康德的自律概念\footnote{《早期神学著作集》，第221页以下。}。与之相反，青年黑格尔的批判意图在于反对康德（同时还有费希特）通过将自由法则和自然规律割裂开来而未能为前者提供充分的范围，因为这种割裂导致实践理性自身饱受“差别”之害。由于黑格尔把康德和费希特哲学中的自由和自然之间的差别关系——黑格尔将这种关系解释为强制和压制的关系——置于批判的中心，并在一个政治整体的理念中寻求将两者结合起来的可能性，这就迫使他回到康德之前的政治哲学传统。

黑格尔在耶拿时期的早期著作中对康德和费希特的批判，就是从这一吊诡开始的。把自由概念扩展为一个政治整体，这一政治整体使个体能够“自由、理想地行动”，这就要求：现在，在伦理理念的建构中，那种遭到压制的“自然”应当重获它的“权利”。《论自然法的科学探讨方式》一文即以此为出发点，这篇论文的标题明白地提到经院哲学的体系形式：它“在实践科学中的地位”\footnote{参见拙文《黑格尔对自然法的批判》（Hegels Kritik des Naturrechts），载《黑格尔研究》第4卷（1967年），第177页以下，本书第42页以下。}。黑格尔在第一部分论述了这门科学的危机。自霍布斯以来，通过对其传统的批判，这一危机就存在了。自然法——同样也包括康德拒绝承认其为实践哲学的政治学和国家经济学——遭受这样的命运，即在现代思想历程中，人们仅将“哲学的哲学成分”归于形而上学，以至于这些科学“完全与理念无缘”\footnote{《全集》第1卷，第437页以下。}。黑格尔认为，对于将建立在经验和技术应用之上的那些科学从形而上学中解放出来，以及形而上学脱离自然、历史中现实物的现实性这样的做法，康德的批判哲学不仅无法逆转，相反，还把它们推向了无路可走的极端。因为批判哲学将自然作为非理性的杂多与作为纯粹统一性的理性对立起来，因而对康德的批判哲学来说，理性——与之一道，还有道德哲学的“纯粹意志”——成为了“一的无本质的抽象”；与此相反，理性则将自然蒸发为“杂多的无本质的抽象”。由此出发，黑格尔返回到了一种历史上先于这种对立的传统形态。这就出现了值得注意的现象：黑格尔试图重新恢复实践哲学中那些从形而上学中解放出来的科学，他不仅接受它们旧的概念形式，而且甚至为作为这些旧的概念形式之基础的“旧的、完全不一贯的经验”辩护。黑格尔在忆及那些将实践哲学的经验置于几何学的认识理想之下，因此必然发现实践哲学的概念形式“既不连贯又相矛盾”的批判家们时说道，一种伟大的、纯粹的直观能够“在其展现的纯粹建筑（在这个建筑中，必然性的联系和形式的统治并不是显而易见的）中将真正的伦理事物表达出来，各个部分以及那些自我修正的规定性的排列组合暗示了那种尽管看不见然而却是内在的、合理的精神”\footnote{《全集》第1卷，第453页以下。}。

在这种直观中，我们很容易发现亚里士多德实践哲学的自由建筑术的特征。通过返回到这种直观，黑格尔对现代自然法的原则和古典政治学的原则进行对照。个体的存在最清晰地表现了两者在出发点上的差异。在城邦伦理中，个体既不执着于自身，也不是通过其自然的特殊性（社交冲动、对死亡的恐惧、自我保存等）与他人结合为一个社会；他不是以抽象的方式——通过契约和个体意志的服从，而是自然地（$\phi\acute{\upsilon}\sigma\varepsilon\iota$），或者如黑格尔所言，以“活生生的”方式与伦理整体合而为一。与将“个体的存在设定为首位和最高原则”的现代自然法相对，古典政治学认为“个别性自身纯属虚无，它与绝对的伦理尊严纯为一物——唯有这种真正的、生动的、并非服从的同一才是个人的真正伦理”\footnote{同上书，第452页。}。我们可以在黑格尔早期实践哲学构想中观察到这种接近古典政治学概念的倾向。在他对伦理学的评价中，这种倾向也表露无遗。尽管伦理学作为德性论，必须是个人的伦理，但是“绝对”伦理也必须在个体上面表现出来，个体完全是否定性的东西，居于否定的形式之下；不仅现代主体性的反思性道德，而且外在地进行自我保证的、被嵌入伦理整体中的个体德性（黑格尔按照古典模型对其进行了自然描述\footnote{同上书，第513页。}）也是“否定的伦理”。这构成了“道德学”与自然法之间的形式区分，然而这种区分并非如康德和费希特所理解的那样，“仿佛它们是分离的，似乎道德学已经从自然法中排除出去了；毋宁说，道德学的内容完全就在自然法之中”\footnote{参见《伦理体系》，《政治学与法哲学文集》，第465页；《全集》第1卷，第510页。}。

这个根本而言为古典的自然法概念及其在实践哲学中“地位”的困境在于，它导致了个体存在以及作为消解和运动环节的否定物的消失。这种困境导致黑格尔无法对现代自然法做出合乎历史的恰当理解，而当黑格尔把亚里士多德关于劳动和行动的理论作为其自然法构建要素接受下来时，这种困境延续了下来。按照黑格尔的理论，伦理的保存要归功于那些个体的德行，个体的活动并不指向一件“产品”——亚里士多德的作品（$\acute{\epsilon}\rho\gamma o\nu$），而是直接在其规定性中摧毁之；它是“没有目的，没有需求，也不涉及与实践情感的关系，没有主观性，没有涉及占有和取得的主观性；相反，此种主观性的目的和产品与其自身一道终止”\footnote{《政治学与法哲学文集》，第467页，另见第463、466页。正是在这里，“步枪”作为一种普遍的、无人身的勇敢和“无人身的死亡”的特别现代的（“机械的”）发明，直接出现了（第467页以下）。}。德性的行动是“政治家和战士”的“绝对无差异的行动”，在《伦理体系》中，黑格尔想将其划给一个“绝对的等级”，我们很容易重新认出这就是古典意义上的市民等级、生而自由的、英勇善战的男性等级。与之相对的是从事“差异”事务（与自然关系相关的“经济”事务）的“正直等级”，这一等级从事“需要的劳动，占有、取得和财产”。这就是现代意义上的不参与战争、局限于私人生活的布尔乔亚（bourgeois）的市民等级。在黑格尔看来，随着奴隶关系的消失，产生了这一等级：“这里，对占有和特殊性的这种沉迷，不再受役于（‘绝对等级’的——笔者）绝对无差异。它尽可能地无差异化了，或者说形式的无差异，个人的存在（Personsein）在人民中反映出来，占有者通过其差异，没有没入整个共同体之中，因此没有陷入人身依附中；相反，其否定的无差异被思为某种实在之物，因此他是市民，是布尔乔亚，并且作为普遍物而得到承认。”\footnote{《政治学与法哲学文集》，第473页。}

另一方面，黑格尔又并不想像亚当·斯密在《国富论》（1776年）中所做的那样，去颠倒行动（Praxis）与制造（Poiesis）的传统关系\footnote{参见《国富论》（The Wealth of Nations）第1卷，第10章；第2卷，第3章。}。作为新兴的国民经济学理论家，斯密勾画了摆脱封建现代剥锁的现代市民社会的模式。在他这里，在古典政治看来保存了城邦的政治家、法律人与战士的行动，被降格为一种非生产性的活动。这种活动没有价值，因为它并没有在一件“作品”中，也就是在一个持续的对象或是一种商品中将自己实现出来。相反，对黑格尔来说，劳动和行动的关系本身不仅仍然有效，而且在其第一个实践哲学体系草案中，构成了他吸收现代国民经济学以及进行等级划分的出发点。这些草稿的大部分矛盾和晦涩可以归结为这种情况，即黑格尔将古典政治图式作为他分析国民经济学的指南。一方面，他将古典政治图式作为范本加以接受；另一方面，他又不得不对它进行修改和补充。而从古典政治图式本身来看，这些修改和补充并不就是一目了然的。

这些矛盾最为明显地出现在《论自然法的科学探讨方式》（1802—1803年）这篇论文中。这里，通过援引柏拉图和亚里士多德\footnote{参见《全集》第1卷，第495页以下。}，《伦理体系》中的“绝对等级”表现为“自由人等级”。不同于“非自由人等级”，这一等级天生就是独立的或者自为的，而非自由人等级则“按其天性不属于自己，而属于他人”\footnote{译自亚里士多德：《政治学》第1卷，第5章，1254 b 20-23；此外，参见同书第1卷，第4章，1253 b 27-30；《形而上学》第1卷，第2章，982 b 25-26。}。黑格尔继续论述道，第一等级的“劳动”并不限于个别的特定产品，而是指向“伦理有机体之整体的存在与维持”，并将自己展现为整体的“鲜活的运动和……神圣的自我享受”\footnote{《全集》第1卷，第494页以下。黑格尔直接提到了柏拉图和亚里士多德的政治理论：“亚里士多德为这个等级规定的事务，就是希腊人所言的‘$\pi o \lambda \iota \tau \epsilon \acute{\upsilon} \epsilon \sigma \theta \alpha \iota$’，在人民中生活，同人民生活，为人民生活，过一种普遍的、完全属于公共事务的生活，或是从事哲学思考。这两种事务，柏拉图以他更高的生命力认为它们是无法分开的，是完全连在一起的”（第495页）。}。因此，黑格尔赞同传统的实践优先的理论，实践的目的就是纯粹的运动（$\kappa\acute{\iota}\nu\eta\sigma\iota\varsigma$）和活动（$\acute{\epsilon}\nu\acute{\epsilon}\rho\gamma\epsilon\iota\alpha$），而不是一种产品的规定性\footnote{参见《形而上学》第9卷，第6章，1048 b 18-36；第8章，1050 a 21-31。}。精神概念最初在这个图式中得到了实现，并且也正好将表达了伦理行动的自我返回的（“内在的”）运动表达了出来。他在1803—1804年耶拿讲座中讲道：“人民的精神必须自身永恒地成为作品。或者说它只是作为一种朝向精神的永恒生成。”\footnote{《耶拿实在哲学》第1卷，J. 荷夫迈斯特出版，莱比锡，1932年，第232页。}然而，对黑格尔来说，在一族人民中，那种每次从自身中产生出一个产品的活动——现代意义上的与自然打交道的劳动（在这种交道中，运动超出自身，凝结于作品中）——也属于精神的这种“在作品中存在”（Am-Werke-Sein）。因此，黑格尔也在自我组织的伦理之内为“为了个人而劳动”的非自由人等级指派了一个（尽管是从属性的）地位。在黑格尔看来，自我组织的伦理必然被“划分为一个绝对接纳无差异之中的部分（公民的实践，其作品为维系城邦的政治活动自身——笔者），以及一个实在之物本身在其中是存在着的，因而是相对同一的部分（奴隶、手工艺人等生产出产品的制造活动[$\pi o\acute{\iota}\eta\sigma\iota\varsigma$]），这个部分只在自身中带有绝对伦理的反映”\footnote{《全集》第1卷，第494页。}。

不同于现代自然法理论的体系基础，在黑格尔的思想中是政治经济学体系的劳动概念提供了相对的辩护。这种辩护，我们在自然法论文中即已见到，在实践哲学的新建筑术中也同样清晰可见。一方面，黑格尔希望通过一种“伦理自然”法的建构，消除自然法与政治学之间的近代分裂以及私法范畴（比如说契约）的主导地位——特别是在国家法和国际法中，这些私法概念似乎已经变成了“自然法的家庭内部事务”\footnote{同上书，第410页以下。}。这样，在黑格尔试图建立起来的实践哲学体系中，自然法又重新与政治学和伦理学一道，落入根本而言需要从政治出发得到理解的“伦理”概念之下\footnote{同上书，第509-513页。}。同时，对黑格尔而言，政治经济学也是一门非常重要的“实践”科学。由此，它也就取代了旧的家政学，显现为体系的基础。其次，接踵而至的是与之相关的“形式”法。于此，我们主要看到的是旧的、处理私法材料的自然法，最后便是与前两者分离开来的、因成为“绝对”而摆脱了一切经济学和形式之“纠缠”的伦理：“这个领域的概念，作为实在的实践[领域]（Praktische），主观地看，就是感觉或者身体的需要和享受，客观地看，就是劳动和占有；这个实践[领域]（如按照其概念发生）被接纳到无差异当中，就是形式的统一性，或者法权（Recht），法权在这个实践[领域]中是可能的。第三个东西作为绝对物或者伦理，超乎二者之上。”\footnote{《全集》第1卷，第494页。}

毫无疑问，黑格尔由于采取了这样的一种划分方式，既居于传统的中心，又与传统相冲突。与18世纪对实践哲学的材料进行重组的企图相比，黑格尔将体系划分为政治经济学、法（即自然法）与伦理（即政治学和伦理学）的做法已经在自身中固定下来。而这又与自然法论文中政治经济学体系所发现的独特的承认的“等级”形式有关。就此而言，这种体系安排最清晰地表达了黑格尔建构其首个伦理体系的鲜明意图：将“绝对伦理的理念”与一个“贵族和自由人等级”的伦理现象混同起来，并将劳动等级的“相对伦理”与之对立起来。

黑格尔通过把劳动的领域从前述限制中解放出来，而超出了这种安排。这种超出发生在《伦理体系》（1802年）第一次阐述之后（此次阐述在许多方面都不充分），以及在1803—1804年和1805—1806年的“实在哲学讲座”中。在这两次讲座中，劳动与语言和行动（承认）一起，作为中心环节出现在精神的构成中。精神，并非如自然法论文所假定的那样\footnote{同上书，第509页以下；另见《信仰与知识》，同前，第424页。}，等同于“伦理自然”。同时，“伦理自然”也并不等同于一族“人民”的身位（在这一身位中，精神反思它自身，并且在理智直观的媒介中将自然带回到自身）；相反，取决于对这种带回[自然]本身进行反思，以便使精神的“教化”过程成为主题。因为黑格尔认识到，这种教化过程也同与自然打交道连在一起，伦理行动并不能穷尽其含义，由此，对他来说，“劳动”也就成为了实践哲学的问题。黑格尔勾画了一种劳动理论的基本特征，由此便将其自身的重要性给予了亚里士多德以“制造学”（Poietik）之名进行处理的那个哲学部门。自亚里士多德以来，对这个哲学部门的探讨便没有任何进步，在现代自然法理论（霍布斯、卢梭和康德）中，它也没有取得任何进展。新的制造学是把国民经济学与先验观念论联系在一起的一种结果，这种联系在青年黑格尔的思想道路中就已经有所准备，并在耶拿讲座中得到了实现。劳动本身构成精神的一种基本形式，作为这样的形式，它在实践哲学的体系中不再被贬低为一个次要的环节。

与哲学传统相对立，黑格尔既没有把实践行为局限于人与人之间行动的概念，也没有像康德和费希特那样将其局限于与自身感性相对立的道德主体性的内在活动中，而是将其扩展到人与自然或者（用先验哲学的普遍化术语来说）——主体与客体、自我与非我打交道的维度中。对劳动概念的这种新的解释导致实践哲学的原则和构造形式发生了根本性的转变。在其第一个实践哲学建构中，黑格尔通过个体的[伦理]自然法存在，以及通过将“个体的”劳动以等级的方式排除在伦理理念之外，消解了康德哲学的开端；其新的实践哲学则与他最初的实践哲学建构完全对立，以《纯粹理性批判》在统觉的原始综合统一的思想中阐发的先验自我的概念、以必然能够伴随我的一切表象的“我思”为出发点。然而，黑格尔并不是以康德的眼光，而是以——将这种伴随的自发性及其造成可能经验的统一性的成就解释为先验的“创造”的——费希特的眼光，来看待我思的功能。在从18世纪向19世纪过渡之际，费希特的这种解释使哲学卷入了一场撤除迄今为止各门科学之间所有教条主义樊篱的运动之中。青年黑格尔追随费希特和谢林，最先在耶拿讲座中制定出这门新的作为这场运动之表达的“意识的”先验“历史”的科学。黑格尔与费希特一样，执守先验哲学的“主体”的行为及其产物。这些产物规定了存在的意义，并且不再规定那些自笛卡尔以来即导致了那种没完没了的、在康德学派也还没有得到解决的争端的实体本体论的前提：质料或者形式、物或者自我、存在或者自我意识、主体或者客体是否拥有独立的持存。这些互相排斥的规定尽管一方面从古典的亚里士多德本体论的毁灭中产生出来，但是它们却无法逃脱作为一种独立或不独立、能动的或者被动的存在者的前提的实体本体论原则的魅力。对黑格尔来说，此类规定都是形而上学的抽象物，它们不仅错失了精神的教化过程，而且也错失了作为这些对立面之统一的精神的概念。诚如黑格尔所言，对于这样一种“非理性的争论”存在应该建立在意识的“被动性”还是“能动性”之中，“真正说来没有什么理性的东西可言”。实在论者主张形式、意识的被动性，观念论者把这种被动性归之于质料、存在。然而事实上，“真正说来”，这一争论的中心就是“中项的在自身中进行争论的潜能阶次（Potenz），其中诸规定性自身与关联项（Bezogensein）两者既设定为一又设定为有区别的东西”\footnote{《耶拿实在哲学》第1卷，第214页以下。}。黑格尔不是在其对立中思考对立的两个方面，而是将其思为诸对立的统一，但是他是在先验意识的基础上思考这种统一的。先验意识在其自身中融解了形而上学的僵化谓词，使之流动起来。因此，这种自我意识不仅仅是那种“伴随”表象的运动的“我思”的自然统一性，而是自身即处于运动之中的。黑格尔将这种运动表现为自我的历史，表现为自我与他者在语言、工具和作品这些中项之中的中介活动\footnote{同上书，第213页。}。在黑格尔这里，劳动——主体同对象打交道以及主体在作品中的对象化——也就成为意识的中介活动史的一个环节，而作品——作为工具和占有——则成为形式与质料、主体与客体、能动与被动等对立之间的中项。在这一中项中，意识（其自身就是打交道过程中的中项）自我组织起来，并且自我“教化”（bildet）。劳动和教化彼此共属一体。

黑格尔由此颠倒了旧实践哲学的基础。他解决了隐藏在劳动和行动关系规定中的困境。这一困境可以这样说明：如果活动者的活动（$\acute{\epsilon}\nu\acute{\epsilon}\rho\gamma\epsilon\iota\alpha$）为其自身的确证（也就是在伦理—政治实践的情形下）的话，那么这种活动就不能对象化；另一方面，如果活动对象化了的话，那么它就不在活动者当中，而是在其作品中。在作品中，活动同时就已经熄灭了。按照制造（poiesis）的古典模式，在劳动与劳动者之间并不存在任何反思性联系，也没有黑格尔的教化概念所表达的劳动者与作品、主体与客体之间的那种结合和摆动。原因在于活动本身的等级秩序，依据这一等级秩序，行动要高于劳动。实践的明智行动更加完善，因为它在实现出来时将行动者的完善表现了出来，而劳动者则似乎将形式沦丧于质料之中，因此也就不能导向劳动者的完善：“不是制造者的完善，而是作品的完善（non est perfectio facientis, sed facti）”\footnote{托马斯·阿奎那：《神学大全》第I卷第2册，第57题，第5节。}。在技术—实践的意义上与自然打交道时，制造（poiesis）生产出来的东西，并不是服务于作为生产者的生产者——比如说，使他在行动的过程中认识人类的事务——而是服务于这种实践对制造之物的使用。前工业社会的模式就是建立在这种使用之上，这就意味着，建立在那种按照定义剥夺了生产者的明智（phronesis），也就是实践智慧的能力之上。黑格尔通过他对国民经济学劳动过程的分析摧毁了这一模式。他的新模式取代了旧模式，它建立在经过先验哲学改变了的本体论基础之上，黑格尔的外化范畴则凝聚了先验哲学的这个本体论基础。劳动的过程无非就是进行制造，同自然打交道，对外物进行改造。然而黑格尔并不将这一过程理解为把一种外在于劳动者本人的形式转加到一种外在的质料之上，而是理解为劳动者的外化。而且，他还以双重化的方式进行理解：(1) 作为欲望的外化。欲望是“内在地扎根”于个体之中的东西，就此而言，也就是一种普遍的东西；(2) 作为同样方式的“内在”形式的外化，也就是意识的外化。黑格尔把劳动定义为“意识把自身造就为物”（Sich-zum-Dinge-Machen des Bewußtseins）\footnote{《耶拿实在哲学》第2卷，第214页。}。在劳动中，意识失掉了其主观的形式，以便在更高的层次上（在作品的客体化中）获得它，也就是说直观它。由于主体通过劳动而适应自然的联系，因此劳动过程就将主体外化为物；同时，因为主体也在加工的对象中使自己成为对象，也就从外化中摆脱出来。这种双重否定的运动不仅消除了质料和形式的简单对立，而且消除了在古典模式结构中发挥了重要作用的使用的优先性。概而言之，不同于亚里士多德以及追随他的前工业的制造学传统（技术理论或者工艺学），黑格尔并不是从劳动过程的终点而是从它的起点出发，对它进行解释：劳动源于欲望的否定物，在劳动的进展中，它一直都盯着欲望，同时也坚持要否定其所加工的对象：劳动是“受阻的欲望”，并且作为这样的东西而施与教化——否定之否定。

由此出发，黑格尔在《精神现象学》的那个著名章节阐述了主奴辩证法。在一定程度上，主奴辩证法也可视为劳动和行动的辩证法。从内容的阐述来看，这一章紧接在城邦伦理的结构之后。在城邦伦理中，相对于劳动等级及其陷入存在者的个别性之中的沉沦存在，“自由人等级”基于其“赴死的能力”——以及战士的实践——而获得了绝对的优先性。现在，黑格尔认为个别性环节已经在劳动中得到了扬弃，并且为该环节在对象中的持存进行了辩护，这就表明，黑格尔在此期间究竟在多大程度上超出了这种立场。在黑格尔看来，正因为主人对奴隶的统治与一种自身并不产生任何产品的行动相连，因此它就必须瓦解。尽管如此，古典意义上政治行动的环节也没有在此章中再度出现。因此，物品生活（Gut-Leben）不是解释为使用，而是解释为享受，由此，黑格尔就把他自己模式的出发点——欲望与劳动的联系——回植到亚里士多德的模式之中。这样一来，依照黑格尔的观点，在主人这一方面，那个按照古典模式为行动比劳动优越奠基的环节就具有一种缺陷：不能在一个产品中对象化。对于辩证的分析来说，这种观点证明它自己只是假象。因为物在劳动中保持了它的独立性，因此“对于物的非本质的联系方面”便落到了奴隶上面。对于主人来说，对象在享受中消失了；而对于奴隶来说，对于对象的否定联系便成为了“对象的形式并且成为了一种有持久性的东西，这正因为对象对于那劳动者来说是有独立性的”\footnote{《全集》第2卷，第156页。（中译文引自〔德〕黑格尔：《精神现象学》上卷，贺麟、王玖兴译，商务印书馆1979年版，第130页。——译者）}。劳动并没有消逝在加工的对象中，而是构成了欲望和对象之间的持久运动，“否定的中项”或者“赋形的行为”，对劳动者来说，这种行为就像古典观点中行动对于行动者一样内在。“这个否定的中项或赋形的行为同时就是个别性，或者意识的纯粹自为存在，意识现在在外在于它的劳动中进入持久的要素之中；由此，那劳动着的意识就直观到那个独立的存在就是它自身。”以及“由于形式被设定在外，所以对服务意识来说，形式并不是一个有别于它的东西；因为这形式正好就是它的纯粹的自为存在。在形式中，它的这个自为存在成为了真理。”因此，正是在似乎只具有异己意义的劳动中，服务的意识通过重新发现自己、通过自身获得了本己的意义\footnote{同上书，第157页。[这两段引文，贺译本译为“这个否定的中介过程或陶冶的行动同时就是意识的个别性或意识的纯粹自为存在，这种意识现在在劳动中外在化自己，进入到持久的状态。因此那劳动着的意识便达到了以独立存在为自己本身的直观。”以及“奴隶据以陶冶事物的形式由于是客观地被建立起来的，因而对他并不是一个外在的东西而即是它自身；因为这形式正是它的纯粹的自为存在，不过这个自为存在在陶冶事物的过程中才得到了实现。因此正是在劳动里（虽说在劳动里似乎仅仅体现异己者的意向），奴隶通过自己再重新发现自己的过程，才意识到他自己固有的意向。”见《精神现象学》上卷，第130、131页。——译者]}。用一句简短的拉丁语来说就是：factio perfectio facientis est（制造是制造者的完善）\footnote{参见罗森克兰茨：《黑格尔传》，第133页，另见《全集》第1卷，第495页以下；《政治学与法哲学文集》，第467页。}。

\section*{3}

如果我们想要理解黑格尔哲学中的客观精神的概念和安排，就必须要对劳动和行动之间关系的变化予以解释。正因为黑格尔在吸收实践哲学的过程中，不得不修正实践哲学的基础，因此他与传统的争论就使他远远超出了传统以外。这一点可以从1805—1806年的耶拿讲座中得到具体证明。在某些方面，这一讲座放弃了之前的“伦理自然”建构，同时也就舍弃了相应的所有实践哲学的划分尝试。在1801年到1804年之间，精神哲学根本而言与伦理理论一致，还没有进行内部划分。精神首先是一族人民的“绝对”精神，它的其余环节不是对民族精神开端的补充，就是对其终点的扩展。在这一时期，甚至宗教和哲学都属于伦理。依据柏拉图-亚里士多德模式，黑格尔将哲学比作一族人民的“绝对劳动”：正如牺牲的伦理行动的纯粹活动证明自己是为了政治整体的，哲学在对“个别性”的完全否定中将人民精神提升为“自然世界和伦理世界的精神”\footnote{《政治学与法哲学文集》，第465页。}。另一方面，在“伦理精神”和“绝对精神”同一的前提下，也就不可能对个体的（主观）精神进行独立的探讨。因此，在黑格尔看来，个体“以一种永恒的方式”直接存在于伦理整体之中，“其经验的存在和行为是完全普遍的存在和行为；因为个体并不是行动着的个体物，而是寓于他之中的普遍的绝对精神”\footnote{参见《耶拿实在哲学》第1卷，第230页；第2卷，第210页以下。}。

这里，对劳动概念的全新阐述导致了精神哲学体系安排上的深刻转变。此时，1803—1804年讲座所阐发的精神教化的思想，也就是“绝对伦理”从在自我行动的自我中被给予的意识阶段（“理智”和“意志”）“发展”出来的思想，开始发挥出全部效力。由此，也就抹平了自然法论文和《伦理体系》在“自然伦理”或者相对伦理（经济学）与“绝对伦理”（政治学，它对“政治”经济学进行不断的限制）之间划下的鸿沟。根本而言，有两个理由：(1) 借助国民经济学的劳动概念，“个体性”环节在与自由和平等的个体相互承认有关的行动过程中得到了辩护\footnote{《耶拿实在哲学》第2卷，第210页。}；(2) 在此之前，黑格尔将法权概念作为经济-实践领域的单纯“形式”统一性来处理，“绝对物或伦理物”毫不相干地漂浮于这种统一性之上。与此不同，现在黑格尔把法权概念作为“一般伦理”的环节从“承认的运动”中产生出来，同时，通过“承认的运动”，之前分离开来的体系各部分也得以联系起来。尽管如此，黑格尔并没有按照习惯的方式接受自然法的这个核心部分（作为使得从自然状态过渡到“公民状态”得以可能的契约论建构），而是对自然法范畴和政治经济学范畴进行对照和比较。他发现了对他自己和自然法理论家来说仍然隐藏着的东西：法权人格的概念，以及从“个体意志”到“普遍意志”的让渡概念中所包含的个人的自由和平等，是以通过劳动从自然中解放出来为前提的。黑格尔认为卢梭和康德的普遍意志不过是一种抽象物，因为它没有作品，其自身缺乏客观的联系和持久的存在。尽管承认的运动指向的是他人的自我认识而非一物\footnote{同上书，第113页。}，但是得到承认的意志的内容是经过了劳动和占有的中介作用的。只有这样，普遍的意志才是“现实的精神”。在其中，意志和理智如同劳动和行动一样彼此联系起来。最初包含在普遍意志概念之中的东西，并不是像意志的联合和契约这样的社会和历史的抽象构造，而是从人和自然的中介过程中产生出来的那些环节：劳动、产品和占有。这些环节先于所有的意志关系：“精神既不是作为理智也不是作为意志，而是作为本身就是理智的意志，才是现实的……在这种本身就是理智的意志中，抽象的意志必须自我扬弃，或者作为被扬弃了的，在承认（Anerkanntsein）的要素中、在这种精神的现实性中产生出来。就像之前劳动转变为普遍的劳动一样，占有也由此转变为法权；之前作为家庭财产的东西……成为了所有人的普遍作品和享受。”\footnote{《耶拿实在哲学》第2卷，第263页以下。}只有通过关联到一件作品，在这件作品中，普遍意志的概念、理智和意志的统一（“自我”）给予自身以定在，[只有这样，]普遍意志才是“现实的精神”。这种“现实的精神”还不是1817年的《哲学全书》引入的“客观精神”，但也是迈向客观精神的一步。黑格尔通过区分作为定在的教化的外化（这种外化同时就是对象化或客观化）和定在着的世界本身的外化完成了这一步。这一区分为1805—1806年讲座的结论奠定了基础。这个讲座的结尾不再是把伦理提升到宗教，而是把艺术、宗教和科学作为“绝对自由精神”的特殊阶段。绝对自由精神把迄今为止的所有规定都收回到自身，并且产生出“另一个世界”\footnote{参见上书，第212页，附释2。参见哈贝马斯：“劳动与相互作用——黑格尔耶拿精神哲学评注”（Arbeit und Interaktion. Bemerkungen zu Hegels Jenenser Philosophie des Geistes），载《自然与历史：卡尔·洛维特70大寿文集》（Natur und Geschichte: Karl Löwith zum 70. Geburtstag），斯图加特/柏林/科隆/美因茨，1967年，第132页以下、154页注释3。然而哈贝马斯忽略了耶拿精神哲学的现象学倾向。}。这里，一方面，差异似乎已经达到了“绝对精神”，另一方面，“现实的精神”和“主观精神”之间却没有清楚的界分。黑格尔之所以还没有认识到这一术语，与他在耶拿讲座草稿中对精神哲学的现象学倾向有着最为紧密的联系。尽管耶拿讲座从包含了认识（理论自我）和意志（实践自我）的意识概念中发展出了“现实的精神”，但由于现象学方式所要求的对象的相关性，后来所言的“主观精神”和“客观精神”的要素就在这里不断地混杂在一起：想象力、回忆与符号分别同自由意志、工具与技术的狡计一一对应。因此，在理论自我和实践自我之间也就并不存在真正的综合\footnote{《耶拿实在哲学》第2卷，第199页以下。}；实践自我的发展以极不清晰的方式汇入家庭之中\footnote{《哲学入门》（Philosophische Propädeutik）第3篇，第2章，第3部分，第128-129节、第173节以下（《黑格尔全集》历史考订版中对黑格尔纽伦堡时期这一精神学说的划分与里德尔此处引述的荷夫迈斯特的划分不尽相同。前者的划分如下：1. 精神的概念。A. 理论精神。B. 实践精神；第二章“实在精神”；第三章“在其纯粹阐述中的精神”。见G. W. F. Hegel, Gesammelte Werke. Band 10.1. Felix Meiner Verlag Hamburg 2006, S. 342, 353, 359, 362。中译本参见《黑格尔全集》第10卷，张东辉、户晓辉译，商务印书馆2012年版，第273、286、292、295页。——译者）。参见法尔肯贝格（R. Falckenberg）：《黑格尔客观精神的实在性·哲学及其历史论稿》（Die Realität des objektiven Geistes bei Hegel. Abhandlungen zur Philosophie und ihrer Geschichte）第25册，莱比锡，1916年，第62页。即便我们认为罗森克兰茨并没有认识到这些联系，并且人为地破坏了文本，但是这种模糊性依然存在。参见F. 尼科林：《黑格尔的主观精神理论著作》（Hegels Arbeiten zur Theorie des subjektiven Geistes），见《认识与责任》（Erkenntnis und Verantwortung），杜塞尔多夫，1960年，第367页（Th. 利特纪念文集）。}，以便再次开始新的为了承认的斗争。

同样的模糊性重又出现在《纽伦堡哲学入门》中。这里，“精神科学”分为三节：(1) 在其概念中的精神，(2) 实践精神，以及 (3) 在其纯粹阐述中的精神。第一节处理作为（理论的）理智的精神，并以“理性思想”结束（第170-172节）。它并没有实现与“实践精神”的综合，相反，实践精神单独包含了整个伦理学说，由此，伦理学说便以吊诡的方式从属于心理学体系之下。直到海德堡《哲学全书》（1817年）的概念构成才终结了这种追随18世纪90年代对个体灵魂学说的早期划分尝试\footnote{参见罗森茨威格：《黑格尔与国家》（Hegel und der Staat）第2卷，1920年，第83页以下。}而来的摇摆倾向。这部《哲学全书》产生了那种综合，并将“客观精神”定义为“理论精神和实践精神的统一”\footnote{《哲学全书》，第400节。《全集》第6卷，第281页。}。局限于意识之“自我”的现象学限制考察被取消了。“精神”是超越于思维、意志和对象之上的统一，唯有《逻辑学》的进程才能对这种统一进行阐述。由此造成段落安排上的根本差异：1807年《现象学》中的精神：精神的第一个阶段（尚未称作“客观精神”）是把“个体性”环节从自身中排除了出去的古代城邦的实体性伦理，“客观精神”的第一个阶段是“人格”的抽象法。在其中，通过作为普遍和个别之统一\footnote{《逻辑学》，第二部分，拉松出版，莱比锡，1951年，第219页以下。}的逻辑概念自身，它[指抽象法——译者]得到了绝对的辩护。“现象学的名称”与客观精神的意义的一致之处在于，精神在这里也意味着那种“普遍的作品”（而不仅仅是思维与意志的统一、“普遍意志”），“它通过所有人和每一个人的行动，作为他们的统一性和同一性而创造出来”\footnote{《全集》第2卷，第336页。（贺译本译为：“一切个人和每一个人通过他们的行动而创造出来作为他们的同一性和统一性的那种普遍业绩或作品。”见《精神现象学》下卷，贺麟、王玖兴译，商务印书馆1979年版，第2页。——译者）}。但是精神并不是一个自由意志的“客观化”的统一性概念，而是分化为意识的各个不同的“世界”。它在古代城邦中具有实体性，在现代教化王国中有外在的定在，而最终在普遍意志和道德中拥有知识和意志的自我确定性，但没有在一个积极的作品中得到实现。

《哲学全书》对客观精神的处理反其道而行之。它从道德的意志概念开始，这个概念乃是作为自由的起源本身的“普遍”意志，亦即在思维中把握自身的意志。就此而言，这种处理建立在由卢梭、康德和费希特开辟的现代法哲学和国家哲学的基础上。然而，对黑格尔来说，普遍意志既不是客观精神学说的起源，也不是这种学说的目的。之所以不是起源，是因为它是在与自然打交道的过程中获得的；之所以不是目的，是因为它自身必须在一个作品中得到实现。由于黑格尔把自由意志定义为自我规定的且给予其规定以外部实在性的精神\footnote{参见《纽伦堡哲学入门》（Nürnberger Propädeutik），第3篇，第2章，第3部分，第173节。《哲学全书》，第1版，第400节；第3版，第484节。}，所以他把自由意志重新纳入康德和费希特将其从中分离出来的客观定在的领域，由此，这种定在——也就是人的自然和历史（法、国家等）——即表现为由自由意志所“设定”和渗透的。自由只有在这种前提下，即它已经在历史中得到实现并且内在于实定法和国家之中，才能理解为“法的理念”。普遍说来，客观精神的核心范畴，即法，乃是意志在其自由的自我规定中自己给出的定在领域——实现了的自由的王国，精神的世界从它自身中产生出来\footnote{《法哲学原理》第4节。《全集》第7卷，第50页。参见同书第1节，补充，第38页以下。}。

这就是黑格尔的客观精神学说超越其18世纪前辈和传统实践哲学理解之处。自由的、自我客观化的意志的概念在自身中把那些在18世纪前辈和传统实践哲学处依然保持分裂的东西统一了起来：与伦理行动中众多个体意志彼此之间的关系的关联，以及与人类需要本性的关系及其通过劳动与外部自然物之间的中间活动之间的关联\footnote{参见《哲学全书》，第3版，第483节：“自由意志在它身上首先直接具有种种区别。自由是它的内在的规定和目的，并与一个外在的遇上的客观性相联系，这客观性分裂为：具有种种特殊需要的人类学东西，种种为意识而存在的外部自然事物，以及个体意志对个体意志的关系”；另参见《哲学全书》，第2版，第488节到491节。}。行为与活动的这种统一最终表明了精神自身的概念结构\footnote{参见《纽伦堡哲学入门》，第176节。}。尽管作为哲学的原则，精神是绝对的第一者并且构成自然的前提，但黑格尔还是用出自人与自然打交道的领域的那些范畴来解释精神的存在方式。精神不是静止的存在，而是过程、活动和劳动：“但是精神只有扬弃其直接的存在，方才存在。如果精神仅仅存在，那它就不是精神。因为它的存在正在于，作为自为存在着的精神，以自身为中介，自为地存在。”\footnote{《宗教哲学讲演录》（Vorlesung über die Philosophie der Religion）第1卷，G. 拉松出版，莱比锡，1925年，第70页。}在黑格尔这里，“精神的劳动”不是比喻，而是道出了它的本质：成为它自己的产品、成为自我对象化和作品。黑格尔在历史哲学讲座中说，精神“根本而言就是行动，它把自己造就为它自在地所是的东西，造就为它的行迹、它的作品；因此，它成为了自己的对象，使自己作为一个定在出现在自己面前”\footnote{《历史中的理性》（Die Vernunft in der Geschichte），J. 荷夫迈斯特出版，汉堡，1955年，第67页；另参见同书，第55页以下、123、135页。}。

在黑格尔体系内，这种按照劳动和行动相统一的模式对客观精神所做的解释导致了实践哲学的地位发生了根本的改变。在《精神现象学》之后，黑格尔在其体系构建中开始撰写《逻辑学》（1812—1816年），然后出版了作为整体纲要的《哲学全书》（1817年），《法哲学原理》于1821年出版，是他的最后一部著作，而且还是由他本人亲自出版的。这部著作对客观精神学说做出了真正的阐述。这里，黑格尔明确指出了继踵理论哲学与实践哲学的传统体系划分的意图，他在序言中说道：“关于思辨认识的本性，我在我的《逻辑学》中已予详尽阐述。”法哲学不仅建立在逻辑学的辩证方法上，而且还是逻辑学的对应学科，因为它阐述了“实践知识”的本性。黑格尔在柏林时期仍然能够并非毫无理由地按照理论哲学和实践哲学的“习惯划分”模式来解释它们的关系\footnote{参见《纽伦堡著作集》（Nürnberger Schriften），J. 荷夫迈斯特出版，莱比锡，1938年，第XIV页。}。不过，这种划分并不十分契合黑格尔体系的概念理解和结构。对黑格尔而言，逻辑学并不像旧形而上学那样，只是一门哲学的“理论”学科，而是作为“纯粹科学”构成了哲学的第一部分；这一部分要先于体系的其他两个部分，也就是自然哲学和精神哲学，由此，也就取消了逻辑学与自然法的并列\footnote{参见《逻辑学》，序言。《全集》第4卷，第13页。}。因为就事情本身而言，自然法在体系中的位置居于精神哲学当中。在黑格尔这里，“实践哲学”的概念和内容处于“客观精神”的名称之下。为此，体系就必须抛弃旧的“实践哲学”的名称，因为它既不契合于作为现代世界及其社会之基础的劳动，也不契合于精神客观化为其产品的本体论上的必然性。精神的主题是处于法和道德、家庭和社会、国家和历史诸领域之中的人类世界的行为。居于主观精神（人类学、心理学）与绝对精神之间的客观精神学说包含了那些在它产生出来的世界中既有其基础也有其目的之人的行动；此种行为就其意义而言也不是超越的，而是与人类世界及其“客观精神”不可分离的。在黑格尔这里，经院哲学所称的伦理学、家政学和政治学以及与自然法相关的东西，全都落入这个概念之下。由此，一方面设定了同这些科学之间的差别。对黑格尔而言，这些科学，就像旧形而上学、心理学、宇宙学以及自然神学一样，由于哲学的先验转向而“连根拔除了”\footnote{参见《哲学全书》，第1版，海德堡，1817年，第300节（原注似应为第300节，原文标为第299节。——译者）：“对于我们，精神以自然为它的前提，它是自然的真理性。在这种真理性中，自然在精神的概念面前消失了，而精神则作为理念产生了自己，理念的客体同主体一样，都是概念。”（《全集》第6卷，第227页）[中译文参见〔德〕黑格尔：《哲学科学全书纲要》（1817年版），薛华译，北京大学出版社2010年版，第155页。——译者]}。因为在其行动中实现了康德的那个“纯粹意志”的、自由的概念的客观精神，已经克服了自然（人的自然与人的外部世界的自然）的束缚，也就不同于经院哲学传统中的遵循自然规律的自由行为的实现\footnote{《哲学全书》，第3版，第24节，补充2，《全集》第8卷，第87节；《哲学史导论》（Einleitung in die Geschichte der Philosophie），第3版，J. 荷夫迈斯特出版，汉堡，1959年，第110页以下、223页以下；《历史中的理性》，第5版，J. 荷夫迈斯特出版，汉堡，1955年，第54页以下。}。对于黑格尔而言，“客观”精神之为精神，通过实现其“自在”存在着的自由，就不是一成不变的自然规律，而是使自然本身成为一种由精神所设定的、与精神一道得到实现的客体：自由就是在他在中自在地存在（Beisichsein im Anderssein）\footnote{《哲学全书》，第3版，第482节，附释，《全集》第10卷，第380页。（中译文参见〔德〕黑格尔：《哲学科学百科全书III：精神哲学》，杨祖陶译，人民出版社2015年版，第273-274页。——译者）}。对自然束缚的这种摆脱，就为客观精神不仅具有一种历史而是自身就是历史性的，提供了更深的理由。因而，就自由的理念来说，黑格尔能够说，它似乎还从未占领过（或者说尚未占领）非洲和东方的整片土地。这种说法同样适用于历史哲学：“[……] 希腊人和罗马人，柏拉图和亚里士多德，甚至斯多亚派都没有自由的理念；相反，他们只知道，人由于出生（作为雅典、斯巴达或者其他城邦的公民）或由于性格的刚强、教养、哲学……而是现实地自由的。自由的理念是通过基督教来到世上的。”

随着洞见自由理念的历史性（这一洞见对其客观精神学说产生了极深的影响），黑格尔就把传统实践哲学的整座大厦撂到了一边。对于现代思想而言，与青年黑格尔由以出发的那种对传统实践哲学的解构相比，这种洞见的影响要深远得多。19世纪都是从历史哲学出发理解黑格尔的客观精神学说的，由此也就必定忽视了客观精神学说发展起来的来源和背景。今天，哲学又要重新理解这些来源。这同时也就意味着，哲学可以从那种黑格尔彻底摧毁了的传统中获得教益。但是，所有的教益——无论人们怎样理解它，以及人们可以同这些教益走多远——都无法回避与黑格尔对传统的批判性解构之间的冲突。黑格尔的论证不仅尖锐，在许多方面令人信服，而且恰恰集中在经院实践哲学过于明显地与旧欧洲社会统治关系关联在一起的那些方面。问题在于，在当前社会历史世界条件下，实践经验的理论是如何可能的。没有人会说黑格尔给出了充分的回答；然而他的功绩在于，他为他的时代，同时也是我们的时代，提出了这个问题。